% ============================================================
%  Chapter 4 - Experiments and Evaluation
% ============================================================
\renewcommand{\tablename}{Table}
\renewcommand{\figurename}{Figure}

\chapter{Experiments and Evaluation}
\label{chap:experiment}

\section{Research Questions}

This chapter presents the experimental procedure and result analysis of the
reinforcement-learning-based traffic signal control system designed in
Chapter~3. The experiments are structured to answer three research questions.
First, whether the RL algorithms (PPO, DQN, DDQN) outperform classical
controllers (Fixed-time, Webster) on real traffic data, and by what margin.
Second, which of the three reward functions under investigation (queue,
pressure, wait-clip) yields the most effective and stable learning signal across
algorithms and scenarios. Third, how the advantage of RL over the baselines
changes as the network scale increases from Cologne-3 (3 signalised
intersections) to Cologne-8 (8 signalised intersections).

\section{Experimental Setup}

\subsection{Evaluation Data}

All models are evaluated on the real traffic demand from the TAPAS Cologne
dataset, corresponding to the morning peak period 7:00--8:00 (3{,}600
simulated seconds)~\cite{tapascologne}. This evaluation dataset is entirely
separate from the synthetic training traffic: the origin--destination matrix
is kept unchanged from the original dataset, ensuring the authenticity of the
reported results. Each configuration is evaluated with 5 independent random
seeds; results are reported as mean $\pm$ standard deviation across seeds.

\subsection{Compared Methods}

A total of 11 configurations are compared: two classical baselines (Fixed-time
and Webster) and nine RL configurations formed by the full combination of three
algorithms (PPO, DQN, DDQN) with three reward functions (queue, pressure,
wait-clip).

\subsection{Evaluation Metrics}

Performance is measured by two complementary metrics. \textbf{Mean Queue} is
the total number of stopped vehicles across the entire network, averaged over
all decision steps within an episode; a lower value reflects less congestion.
\textbf{Throughput} is the total number of vehicles that complete their
journeys within the episode. An effective controller must simultaneously
maintain low queues and high throughput: a controller that reduces queues by
slowing vehicles before they enter an intersection would be penalised by a
lower throughput.

\section{Training Dynamics}
\label{sec:training}

Figure~\ref{fig:train_c8} presents the median reward curves during training on
Cologne-8, with shaded bands representing the interquartile range across
5 seeds. Across all three reward functions, the three algorithms converge
clearly after approximately 100{,}000 steps. PPO reaches a plateau earliest
and maintains the narrowest IQR across all reward functions. DQN converges
stably but more slowly. DDQN converges but continues to oscillate slightly
after 200{,}000 steps under the queue and pressure rewards. The most notable
case is DDQN with wait-clip on Cologne-8, where the median drops sharply
within the first 50{,}000 steps, recovers gradually, but continues to be
periodically pulled downward with a very wide IQR persisting until 400{,}000
steps, indicating that this configuration produces less stable results across
seeds, consistent with the high standard deviation recorded in the evaluation
results.

Figure~\ref{fig:train_c3} presents the corresponding curves on Cologne-3.
Unlike Cologne-8, none of the algorithms exhibit a clear convergence trend
over the full 500{,}000 steps. The median reward oscillates within a narrow
band close to zero without any sustained improvement. One possible contributing
factor is that fewer intersections means fewer samples are accumulated per
rollout for the shared policy within the same timestep budget, making
convergence more difficult.

\begin{figure}[htbp]
    \centering
    \includegraphics[width=\textwidth]{fig_training_c8.png}
    \caption{Training reward curves on Cologne-8. Solid lines show the median
    reward smoothed over a sliding window; shaded bands represent the
    25th--75th percentile across 5 seeds.}
    \label{fig:train_c8}
\end{figure}

\begin{figure}[htbp]
    \centering
    \includegraphics[width=\textwidth]{fig_training_c3.png}
    \caption{Training reward curves on Cologne-3. Solid lines show the median
    reward smoothed over a sliding window; shaded bands represent the
    25th--75th percentile across 5 seeds.}
    \label{fig:train_c3}
\end{figure}

\section{Evaluation Results}

Tables~\ref{tab:results_c8} and~\ref{tab:results_c3} present the full
evaluation results on Cologne-8 and Cologne-3 respectively.

\begin{table}[htbp]
    \caption{Evaluation results on Cologne-8 (mean $\pm$ std, 5 seeds).
             Lower mean queue and higher throughput are better.}
    \label{tab:results_c8}
    \centering
    \begin{tabular}{lcc}
        \toprule
        \textbf{Method}
            & \textbf{Mean Queue (veh)}
            & \textbf{Throughput (veh)} \\
        \midrule
        Fixed-time       & 16.42             & 2{,}001 \\
        Webster          & 45.85             & 1{,}934 \\
        \midrule
        PPO-queue        & $2.03 \pm 0.05$   & $2{,}014 \pm 1.4$ \\
        PPO-pressure     & $1.93 \pm 0.03$   & $2{,}017 \pm 0.4$ \\
        PPO-wait-clip    & $2.01 \pm 0.04$   & $2{,}014 \pm 1.9$ \\
        \midrule
        DQN-queue        & $2.30 \pm 0.49$   & $2{,}014 \pm 1.1$ \\
        DQN-pressure     & $2.33 \pm 0.46$   & $2{,}015 \pm 1.1$ \\
        DQN-wait-clip    & $2.49 \pm 0.61$   & $2{,}013 \pm 2.5$ \\
        \midrule
        DDQN-queue       & $2.56 \pm 0.67$   & $2{,}013 \pm 2.3$ \\
        DDQN-pressure    & $2.53 \pm 0.39$   & $2{,}015 \pm 1.1$ \\
        DDQN-wait-clip   & $12.51 \pm 23.55$ & $1{,}987 \pm 61$  \\
        \bottomrule
    \end{tabular}
\end{table}

\begin{table}[htbp]
    \caption{Evaluation results on Cologne-3 (mean $\pm$ std, 5 seeds).
             Lower mean queue and higher throughput are better.}
    \label{tab:results_c3}
    \centering
    \begin{tabular}{lcc}
        \toprule
        \textbf{Method}
            & \textbf{Mean Queue (veh)}
            & \textbf{Throughput (veh)} \\
        \midrule
        Fixed-time       & 14.25             & 2{,}797 \\
        Webster          & 33.38             & 2{,}743 \\
        \midrule
        PPO-queue        & $6.10 \pm 9.35$   & $2{,}754 \pm 152$ \\
        PPO-pressure     & $40.20 \pm 26.41$ & $1{,}882 \pm 717$ \\
        PPO-wait-clip    & $4.66 \pm 6.14$   & $2{,}768 \pm 121$ \\
        \midrule
        DQN-queue        & $7.91 \pm 7.22$   & $2{,}713 \pm 147$ \\
        DQN-pressure     & $2.82 \pm 0.32$   & $2{,}819 \pm 0.8$ \\
        DQN-wait-clip    & $33.47 \pm 17.44$ & $2{,}127 \pm 489$ \\
        \midrule
        DDQN-queue       & $4.91 \pm 5.22$   & $2{,}754 \pm 107$ \\
        DDQN-pressure    & $14.27 \pm 11.18$ & $2{,}611 \pm 198$ \\
        DDQN-wait-clip   & $4.25 \pm 5.11$   & $2{,}771 \pm 111$ \\
        \bottomrule
    \end{tabular}
\end{table}

\subsection{Results on Cologne-8}

As shown in Table~\ref{tab:results_c8}, all nine RL configurations outperform
both baselines on both metrics. The best configuration is PPO-pressure, which
achieves the lowest mean queue and the highest throughput among all methods.
Notably, Webster produces the worst results of all evaluated methods, worse
than Fixed-time on both metrics, despite being an adaptive controller.

Among the RL configurations, PPO achieves the highest stability, with a
coefficient of variation below 5\,\% across all configurations. DDQN-wait-clip
is the only configuration that fails severely, with a coefficient of variation
exceeding 100\,\%, consistent with the oscillatory behaviour observed in the
training curves.

\subsection{Results on Cologne-3}

As shown in Table~\ref{tab:results_c3}, the best configuration is DQN-pressure,
which outperforms Fixed-time on both metrics and is the most stable
configuration in the RL group. However, the majority of RL configurations
exhibit large standard deviations, many with a coefficient of variation
exceeding 100\,\%, reducing the reliability of the mean values. PPO-pressure
is the worst configuration in the entire experiment and the only one with a
throughput lower than both Fixed-time and Webster.

\section{Analysis and Discussion}

\subsection{Performance of RL Against Classical Baselines}

On Cologne-8, all RL configurations outperform both Fixed-time and Webster by
a large margin. The improvement over Webster is particularly striking, as
Webster performs worst among all evaluated methods despite being an adaptive
controller. This is because Webster optimises each intersection independently
using a 15-minute moving average and cannot account for inter-intersection flow
dependencies: a long cycle assigned to one node may cause vehicles to accumulate
at adjacent nodes in ways that Webster cannot anticipate. Independent MARL
addresses this by reacting directly to local queue states at every decision step.

On Cologne-3, some RL configurations improve clearly over Fixed-time while
others produce substantially worse results, most notably PPO-pressure. Among
all configurations evaluated on this scenario, DQN-pressure achieves the best
and most stable result, with the lowest standard deviation in the RL group.

\subsection{Effect of Reward Function}

Since the training curves on Cologne-3 do not show clear convergence as
discussed in Section~\ref{sec:training}, the reward function analysis focuses
on Cologne-8, where the models learn stable policies and evaluation results
can be interpreted reliably.

On Cologne-8, the pressure reward yields the lowest mean queue across all three
algorithms, suggesting that the incoming-to-outgoing vehicle imbalance signal
effectively captures congestion dynamics in this multi-intersection urban
network. The queue reward is consistently second-best and produces no failure
cases across any algorithm, making it the most stable reward function across
all three algorithms. Wait-clip performs well with PPO but fails severely with
DDQN, indicating that its effectiveness depends substantially on the algorithm
it is paired with. The wait-clip reward accumulates only after vehicles have
been stationary for several consecutive steps, producing a sparse and delayed
signal, and DDQN-wait-clip records a coefficient of variation exceeding 100\,\%
in evaluation, consistent with the periodic oscillations observed in the
training curves. Overall, pressure is the best choice when combined with PPO
or DQN, while queue is the safest choice when there is no preference for a
specific algorithm.

Regarding the pressure reward on Cologne-3, it exhibits highly divergent
results across algorithms: DQN-pressure achieves the best result on this
scenario while PPO-pressure produces the worst. However, since no algorithm
converges clearly on Cologne-3 as shown in the training curves, these
evaluation results may reflect seed-to-seed variability rather than genuine
policy differences.

\subsection{Effect of Network Scale}

The results show that Independent MARL performs most effectively on Cologne-8
and does not demonstrate consistent advantage on Cologne-3. Cologne-8, with
8 signalised intersections, presents a more complex environment with richer
interactions among nodes, conditions that Independent MARL can exploit through
each agent reacting directly to its local queue state at every decision step,
while the shared policy accumulates experience from all 8 agents in parallel.
On Cologne-3 with only 3 intersections, as discussed in
Section~\ref{sec:training}, fewer agents means sparser learning signal per
rollout, which prevents the shared policy from converging to a stable policy
within the available timestep budget. The results suggest that Independent MARL
achieves its clearest advantage when the number of signalised intersections is
large enough for the shared policy to accumulate sufficient and diverse training
signal.